% Generated by scripts/tikz_reviews.py using the compiled book's labels.
\pdfbookmark[0]{Chapter 4}{comparison-chapter-4}
\ReviewFigure{fig-spline-ortho}
{Orthogonalizing a cardinal spline}
{figures/wavelets/spline-ortho.pdf}
{figures/tikz/wavelets-spline-orthogonalization.pdf}
{Spectral orthogonalization; Proposition \ref{prop-spectral-ortho} and Figure \ref{fig-spline-ortho}.}
{The formula identifies the handwritten example as the orthogonal linear spline. Its triangular generator has compact support, while the orthogonalized function has alternating exponentially decaying tails. The replacement uses the correct Fourier multiplier sqrt(3/(2+cos omega)); its cosine coefficients were evaluated by 8192-point midpoint quadrature. Axis ticks distinguish the generator interval [-1,1] from the displayed orthogonalized interval [-4,4], which uses a compressed horizontal scale. Four small degree examples recover the original spline-family sketches.}
{wavelets-spline-orthogonalization}
{Complete figure}
{4.3}
\ReviewFigure{fig-spline-vs-shannon}
{Sharp and smooth frequency partitions}
{figures/wavelets/spline-vs-shannon.pdf}
{figures/tikz/wavelets-shannon-spline-bands.pdf}
{Multiresolution detail spaces; Figure \ref{fig-spline-vs-shannon}.}
{Interpretation to check: The Shannon row shows the exact low-pass band |omega|<=pi and the next two dyadic wavelet bands. The spline row deliberately shows schematic smooth overlap, not the spectrum of a named spline order. Only a finite frequency window is drawn; the partition identity refers to all scales. Cropped neighboring prose in the scan is omitted. The energy of a scale-j wavelet is multiplied by 2 to the power -j to remove its L2 dilation factor, as the displayed identity explains; this gives the factor two on the outer scale-minus-one band.}
{wavelets-shannon-spline-bands}
{Complete figure}
{4.6}
\ReviewFigure{fig-periodize}
{Wrapping a function around the unit interval}
{figures/wavelets/periodize-1.pdf}
{figures/tikz/wavelets-periodization.pdf}
{On bounded domains; f-periodic is the sum of integer translates.}
{Interpretation to check: The illustrative compact bump is cos squared(2pi(x-0.87)) for |x-0.87|<=0.25, zero elsewhere. On [0,1], only f(x) and its left-shifted copy f(x+1) contribute, so the drawn sum and matching endpoint values are exact. The profile is not presented as a particular orthogonal wavelet; periodization applies to each basis function individually. No unknown profile or empirical values are inferred from the scan.}
{wavelets-periodization}
{Panel 1 of 2}
{4.7}
\ReviewFigure{fig-periodize}
{One representative of each periodic translate}
{figures/wavelets/periodize-2.pdf}
{figures/tikz/wavelets-periodic-translates.pdf}
{On bounded domains; 0<=n<2\textasciicircum{}\{-j\}.}
{Interpretation to check: The selected translation centers lie in the half-open unit interval. The example j=-2 has four distinct centers; the center at one is the same periodic translate as zero. Smooth pulse profiles are schematic and do not assert that this particular shape generates an orthogonal wavelet basis.}
{wavelets-periodic-translates}
{Panel 2 of 2}
{4.7}
\ReviewFigure{fig-haar-inner}
{Haar filter coefficients from overlap integrals}
{figures/wavelets/haar-inner.pdf}
{figures/tikz/wavelets-haar-inner-products.pdf}
{Discrete wavelet transform; equation \eqref{eq-dfn-h-g}.}
{The original inner products are +1 for the scaling function and +1 or -1 for the wavelet. The chapter filter definition adds a factor 1/sqrt(2); the resulting h=(1,1)/sqrt(2) and g=(1,-1)/sqrt(2) are therefore explicitly distinguished from the raw overlap integrals.}
{wavelets-haar-inner-products}
{Complete figure}
{4.9}
\ReviewFigure{fig-filter-constr}
{Complementary low-pass and high-pass energies}
{figures/wavelets/filter-constr.pdf}
{figures/tikz/wavelets-filter-constraints.pdf}
{Wavelet design; low-pass and QMF constraints.}
{The solid and dashed curves are the exact Haar examples |h-hat| squared=2 cos squared(omega/2) and its pi shift=2 sin squared(omega/2). Their sum is two, and the shifted energy equals that of the QMF high-pass filter. The nearby cropped text in the scan is replaced by the complete relevant equality.}
{wavelets-filter-constraints}
{Complete figure}
{4.24}
\ReviewFigure{fig-vanmoments}
{Vanishing moments and flat filter zeros}
{figures/wavelets/vanmoments.pdf}
{figures/tikz/wavelets-vanishing-moments.pdf}
{Wavelet design; vanishing-moment condition \eqref{eq-cond-vm-fourier}.}
{The drawing separates the low-pass zero at pi from the high-pass zero at zero, avoiding the ambiguous overlaid labels in the scan. The curves are magnitudes for a two-moment QMF example; the derivative statement concerns the smooth filter symbols themselves, including their phases. The displayed zeros have order two, while the formula states the general p-moment equivalence.}
{wavelets-vanishing-moments}
{Complete figure}
{4.26}
